% Preface -- Financial Engineering in Practice
% The title below is not final; change it in this one line if you choose another.
\chapter*{Preface}
\addcontentsline{toc}{chapter}{Preface}
\markboth{Preface}{Preface}

\textbf{Asset pricing}\index{asset pricing|textbf} is the study of why financial assets are worth what they are worth. An asset, such as a stock, a bond or an option, gives its owner a right to uncertain payments in the future, and its price today must reflect both when those payments arrive and how risky they are. Asset pricing explains how markets translate future payoffs into present prices, why some assets earn higher average returns than others, and how those returns compensate investors for bearing risk.

\section*{Questions as a way in}

Asset pricing is full of open questions. Does value beat growth? Does momentum work? Are stock returns predictable? Do optimized portfolios beat naive ones? Can machine learning find hidden alpha?

Practitioners act on the answers, products are sold on them, and many answers are taught as if they were settled. Yet in the research literature, many of these questions remain open. Studies that ask the same question often reach different answers, because they use different data, sample periods, methods or standards of evidence, and later work has overturned or narrowed many early findings. These questions therefore make an excellent starting point for learning, because each one rests on a model, a body of research and data, and an ongoing debate.

Each chapter of this book is organized around one question, such as ``do investors earn the returns their stocks earn?'' A question is a good place to start because it gives the methods a purpose: to understand it, you need the theory behind it, the data used to study it, and the reasons researchers have answered it differently. The book reports where the published evidence stands, and introduces those debates. Every chapter follows the same structure:

\begin{enumerate}
  \item \textbf{Foundations.} Learn the core model and the background the chapter needs, in plain language.
  \item \textbf{How it works.} See how the model turns into measurements and methods that researchers use.
  \item \textbf{The question.} Read the chapter's question, stated precisely.
  \item \textbf{The original evidence.} Learn what the founding studies found and how later researchers responded.
  \item \textbf{Lab.} Work with the chapter's methods yourself, in Python, on free data and on simulations.
  \item \textbf{Where the evidence stands.} See a summary of the published findings and whose evidence supports each one.
  \item \textbf{What practitioners and investors can use.} Take away the lessons that matter in practice, and what remains open.
\end{enumerate}

Because the question comes first, you learn each method when you need it. Newey–West errors, false-discovery rates and walk-forward evaluation become tools for answering a question, not abstract techniques. Each lab also asks you to predict a result before you see it, and a prediction that turns out wrong is one you remember.

The evidence does not point the same way for every question. Some early answers have held up, such as the predictability of volatility; others hold only in weaker form, such as value and momentum; and on some, such as how many published anomalies hold up, researchers still disagree.

Repeated across twenty-three chapters, this way of working becomes a habit. That habit, more than any single result, is what I hope you take from this book: before believing a result, ask ``compared to what, out of sample, and after costs?''


\section*{How the book is organized}

The book has nine parts, twenty-three chapters and four appendices. Every chapter is organized around a question, and the chapters are ordered by two rules. Methods climb from simple to complex, so no chapter needs a method taught later. 
Table~\ref{tab:pref-parts} lists the nine parts and, for each chapter, the question it asks and the models and tools it teaches.

\begingroup
\small
\setlength{\LTpre}{6pt}\setlength{\LTpost}{6pt}
\begin{longtable}{@{}r>{\raggedright\arraybackslash}p{0.44\linewidth}>{\raggedright\arraybackslash}p{\dimexpr 0.50\linewidth-4\tabcolsep\relax}@{}}
\caption{The parts and chapters of the book.}
\label{tab:pref-parts} \\
\toprule
Ch. & Question & Models and tools \\
\midrule
\endfirsthead
\multicolumn{3}{@{}l}{\footnotesize\itshape Table~\ref{tab:pref-parts}, continued} \\
\toprule
Ch. & Question & Models and tools \\
\midrule
\endhead
\bottomrule
\endlastfoot
\multicolumn{3}{@{}l}{\textbf{Part I. Returns, luck and data}} \\*
\ref{ch:01} & Does a strong track record prove skill? & Returns, Sharpe ratio, t-statistics, luck by simulation \\*
\ref{ch:02} & Can historical data show what a strategy would have earned? & Adjusted prices; survivorship, look-ahead and micro-cap biases \\
\addlinespace[4pt]
\multicolumn{3}{@{}l}{\textbf{Part II. The classical cross-section}} \\*
\ref{ch:03} & Can size and value explain returns that beta cannot? & OLS, the CAPM, factor models \\*
\ref{ch:04} & Do past winners keep winning? & Portfolio sorts, cross-sectional momentum, crash risk \\
\addlinespace[4pt]
\multicolumn{3}{@{}l}{\textbf{Part III. Predictability through time}} \\*
\ref{ch:05} & Can simple trading rules beat buy-and-hold? & Trading rules, time-series momentum, the Reality Check \\*
\ref{ch:06} & Can valuations predict the equity premium out of sample? & Walk-forward evaluation, out-of-sample $R^2$, forecast tests \\
\addlinespace[4pt]
\multicolumn{3}{@{}l}{\textbf{Part IV. Volatility and risk}} \\*
\ref{ch:07} & How predictable is volatility? & ARCH/GARCH, realized and implied volatility, HAR \\*
\ref{ch:08} & How reliably do VaR models measure risk? & VaR, expected shortfall, backtesting, Basel \\
\addlinespace[4pt]
\multicolumn{3}{@{}l}{\textbf{Part V. Portfolios and dependence}} \\*
\ref{ch:09} & Can optimized portfolios beat an equal split across assets? & Mean--variance, efficient frontier, covariance estimation \\*
\ref{ch:10} & Does diversification fail when you need it most? & Correlation dynamics, tail dependence \\
\addlinespace[4pt]
\multicolumn{3}{@{}l}{\textbf{Part VI. Currencies and interest rates}} \\*
\ref{ch:11} & Do high-interest currencies depreciate enough to offset their yield? & Covered and uncovered interest parity, carry \\*
\ref{ch:12} & What do forward rates say about future interest rates? & Yield curves, forward rates, duration, expectations hypothesis \\
\addlinespace[4pt]
\multicolumn{3}{@{}l}{\textbf{Part VII. The reckoning: from paper to practice}} \\*
\ref{ch:13} & Do investors earn the returns their stocks earn? & Holding-period returns, stopping times, the disposition effect \\*
\ref{ch:14} & Is a profitable backtest profitable to trade? & Spreads, turnover, break-even costs, capacity \\*
\ref{ch:15} & Why might low-risk stocks beat high-risk stocks? & Security market line, Fama--MacBeth, leverage constraints \\*
\ref{ch:16} & How many published anomalies hold up? & Multiple testing, shrinkage, post-publication tests \\
\addlinespace[4pt]
\multicolumn{3}{@{}l}{\textbf{Part VIII. Derivatives and credit}} \\*
\ref{ch:17} & What do credit spreads reflect? & Default, recovery, spreads, the expected-loss spread \\*
\ref{ch:18} & How much option risk does delta hedging remove? & Options, Black--Scholes, the Greeks, dynamic hedging \\*
\ref{ch:19} & Does selling options earn a reliable premium? & Implied volatility, the VIX, the variance risk premium \\
\addlinespace[4pt]
\multicolumn{3}{@{}l}{\textbf{Part IX. Machine learning and foundation models}} \\*
\ref{ch:20} & Can machine learning predict returns better than linear models? & Regularization, trees, neural nets, cross-validation \\*
\ref{ch:21} & Can foundation models forecast better than classical methods? & ARIMA, exponential smoothing, zero-shot foundation models \\*
\ref{ch:22} & Can transformers price stocks better than factor models? & Stochastic discount factor, IPCA, attention \\*
\ref{ch:23} & Can foundation models trained on market data beat GARCH? & Range-based volatility, fine-tuning, synthetic data \\
\end{longtable}
\endgroup

Each question was chosen to do two jobs at once: to matter to practitioners or researchers, and to teach a core model. Every question had to meet five criteria before it earned a chapter.

\begin{enumerate}
  \item \textbf{Anchored in a core model.} Answering it requires you to learn a model that belongs in any financial engineering course.
  \item \textbf{Consequential.} Practitioners act on the answer, products are sold on it, or the question shaped the research literature.
  \item \textbf{Workable with free data.} The central idea of its founding study can be studied hands-on with the French library, FRED and similar public sources.
  \item \textbf{Still debated.} Recent papers disagree, so the chapter can show you live research rather than a settled answer.
  \item \textbf{A spread of outcomes.} In the published evidence, some early answers have largely held up (volatility is predictable), some hold only in weakened form (value, momentum, the option premium), and some have largely failed (forward rates as forecasts, optimized portfolios beating an equal-weighted one). The book does not lean toward debunking.
\end{enumerate}

Working through the questions takes you through the standard financial engineering curriculum. The questions do not cover everything, however. This is not a full treatment of stochastic calculus, exotic derivatives, term-structure models beyond the basics, corporate finance, or market microstructure at the level of the order book.

The book can therefore stand alone or sit beside the classics. Hull, Cochrane, and Campbell, Lo and MacKinlay teach how the models are built and what they imply. This book adds hands-on practice with every method, on data and simulations you run yourself. For readers who do not have those texts at hand, each chapter's Foundations section supplies the background the chapter needs.

Every chapter has a Python lab. The labs are not replications of the studies they draw on, because many of those studies used proprietary or commercial data that most readers cannot obtain, such as the brokerage records behind Odean (1998). Each lab instead borrows the central idea of a study and works with data that anyone can download for free, such as the Kenneth French library, FRED and Goyal and Welch's predictor file, together with simulations in which the right answer is known in advance. The simulations let you see a mechanism clearly, and the real data let you see how it behaves in markets. Where an original study used a paid database, the lab explains how the free substitute differs.

\section*{Foundations sections by chapter}

The Foundations section in every chapter is part of what each chapter expects you to learn.
The table below lists the standard texts that treat its topic in depth. You do not need the books in the Further reading column to follow this one. They are listed for readers who want a deeper treatment of a topic.

\begingroup
\small
\setlength{\LTpre}{6pt}\setlength{\LTpost}{6pt}
\begin{longtable}{@{}r>{\raggedright\arraybackslash}p{0.468\linewidth}>{\raggedright\arraybackslash}p{0.364\linewidth}@{}}
\toprule
Ch. & Foundations covers & Further reading \\
\midrule
\endfirsthead
\toprule
Ch. & Foundations covers & Further reading \\
\midrule
\endhead
\bottomrule
\endlastfoot
\ref{ch:01} & Simple, log and excess returns; the risk-free rate; expected return and volatility; the Sharpe ratio; standard errors and t-statistics & Bodie, Kane and Marcus (2024); Campbell, Lo and MacKinlay (1997) \\
\ref{ch:02} & Corporate actions and adjusted prices; index weighting; CRSP, Compustat, FRED and the French library; delisting; point-in-time data & Bodie, Kane and Marcus (2024); Bali, Engle and Murray (2016) \\
\ref{ch:03} & Describing a distribution; linear regression and the OLS assumptions & Wooldridge (2025); Bodie, Kane and Marcus (2024); Cochrane (2005) \\
\ref{ch:04} & Measuring past performance; cross-sectional and time-series momentum; volatility scaling; underreaction and overreaction & Bodie, Kane and Marcus (2024); Ilmanen (2011) \\
\ref{ch:05} & Moving-average and breakout rules; time-series versus cross-sectional momentum; volatility targeting; rule universes and the SPA test & Campbell, Lo and MacKinlay (1997); Tsay (2010) \\
\ref{ch:06} & Information sets; training, validation and test splits; why in-sample $R^2$ flatters & Hyndman and Athanasopoulos (2021); L\'opez de Prado (2018) \\
\ref{ch:07} & Volatility and variance; historical volatility and EWMA; ARCH and GARCH(1,1); realized and implied volatility & Hull (2022); Tsay (2010) \\
\ref{ch:08} & Value at Risk and expected shortfall; scaling one-day VaR to ten days; when diversification raises VaR & Hull (2022, 2023) \\
\ref{ch:09} & The investor's problem; the efficient frontier; equal weighting (1/N); why the test must be out of sample & Bodie, Kane and Marcus (2024); Luenberger (2014) \\
\ref{ch:10} & Correlation and diversification; why measured correlation rises with volatility; the Forbes--Rigobon correction; tail correlation and tail dependence; the stock--bond correlation & Bodie, Kane and Marcus (2024); Ang (2014) \\
\ref{ch:11} & Spot and forward exchange rates; covered and uncovered interest parity; the carry trade & Hull (2022); Bodie, Kane and Marcus (2024) \\
\ref{ch:12} & Zero-coupon yields; forward rates; holding-period returns; duration; the expectations hypothesis & Hull (2022); Tuckman and Serrat (2022) \\
\ref{ch:14} & The order book; market and limit orders; who provides liquidity; long and short positions; long--short portfolios; turnover & Harris (2003); Bodie, Kane and Marcus (2024); Hull (2022) \\
\ref{ch:15} & Beta, the CAPM and the security market line; alpha & Bodie, Kane and Marcus (2024); Ang (2014) \\
\ref{ch:16} & The efficient market hypothesis; the joint-hypothesis problem; anomalies; replication, reproduction and robustness; publication bias; shrinkage & Bodie, Kane and Marcus (2024); Bali, Engle and Murray (2016) \\
\ref{ch:17} & Credit spreads; default probability and recovery; the expected-loss spread; why spreads can exceed expected losses & Hull (2022) \\
\ref{ch:18} & Payoffs and put--call parity; the Black--Scholes--Merton model; the Greeks; implied volatility; delta hedging and its profit and loss; discrete hedging; the volatility smile & Hull (2022); Shreve (2004) \\
\ref{ch:19} & The VIX; the option seller's payoff & Hull (2022) \\
\ref{ch:20} & Time-series and cross-sectional questions; a simple factor model; why not use every predictor; overfitting and the bias--variance trade-off & James, Witten, Hastie and Tibshirani (2021); Hastie, Tibshirani and Friedman (2009); L\'opez de Prado (2018) \\
\ref{ch:21} & Classical benchmarks; foundation models; contamination; fair evaluation & Hyndman and Athanasopoulos (2021) \\
\ref{ch:22} & The stochastic discount factor and the tangency portfolio; shrinkage SDFs and IPCA; deep SDFs; attention and the transformer & Cochrane (2005); Bali, Engle and Murray (2016) \\
\ref{ch:23} & OHLCV bars; range-based volatility; tokenizing bars; fine-tuning; synthetic market data; the Giacomini--White test & Hull (2022); Tsay (2010) \\
\end{longtable}
\endgroup

Full references for these books are listed alphabetically under Further Reading at the end of the book (page~\pageref{furtherreading}).


\section*{Who this book is for}

This book is for anyone who builds, buys or studies quantitative strategies and wants to understand the methods and the evidence behind them.
Students in financial engineering, quantitative finance and financial data science will use the whole book, labs included, as will quantitative analysts and researchers at asset managers, banks and fintech firms.
PhD students can use it as a hands-on companion to the research literature. Portfolio managers, allocators and serious individual investors can read the Foundations sections and the summaries of the evidence and skim the labs.

You will need basic Python and a working knowledge of introductory statistics, regression and the elements of investments. Everything else is taught along the way.

\section*{How to use this book}

The book covers a wide range of financial engineering.
A one-semester course can follow Parts I to VII and Part IX in order. Instructors can treat Part VIII, on derivatives and credit, as optional. Learning objectives for every chapter, which instructors can use to plan a syllabus or write exam questions, are in the instructor guide in the companion GitHub repository, next to the capstone project.

\section*{An invitation}

Learning by doing is slower at first than learning by reading. You will make predictions that turn out wrong, and you will spend time on ideas that turn out to be harder than they first looked. But you will understand each method because you needed it, and you will remember each result because you worked through it yourself.


\bigskip
\begin{flushright}
\emph{Chris Kuo}\\
\emph{January 2027}
\end{flushright}
